% !TeX root = ../QuestionSet.tex
% ==================== 第9章 概率与统计 ====================
\QuestionSetChapter{概率与统计}

% ===== 9.1 随机事件与概率 =====
\QuestionSetSection{随机事件与概率}

\questionGroup{一、选择题}

% ----------------------------------------
\begin{question}
一个箱子里有5个相同的球，分别以$1\sim5$标号，若每次取一颗，有放回地取三次，记至少取出一次的球的个数$X$，则数学期望$E(X)=$\blank{3cm}．
\end{question}
\begin{answer}{$\dfrac{61}{25}$}
$P(X=1)=\dfrac1{25}$，$P(X=2)=\dfrac{12}{25}$，$P(X=3)=\dfrac{12}{25}$，$E(X)=1\cdot\dfrac1{25}+2\cdot\dfrac{12}{25}+3\cdot\dfrac{12}{25}=\dfrac{61}{25}$。
\end{answer}

% ----------------------------------------
\begin{question}
为研究某疾病与超声波检查结果的关系，从做过超声波检查的人群中随机调查了1000人，得到如下列联表：

\begin{center}
\begin{tabular}{|c|c|c|c|}
\hline
超声波检查结果组别 & 正常 & 不正常 & 合计 \\ \hline
患该疾病 & 20 & 180 & 200 \\ \hline
未患该疾病 & 780 & 20 & 800 \\ \hline
合计 & 800 & 200 & 1000 \\ \hline
\end{tabular}
\end{center}

(1)记超声波检查结果不正常者患该疾病的概率为P，求P的估计值；

(2)根据小概率值$\alpha =0.001$的独立性检验，分析超声波检查结果是否与患该疾病有关.
附$\chi^2=\dfrac{n(a d-b c)^2}{(a+b)(c+d)(a+c)(b+d)}$，

\begin{center}
\begin{tabular}{|c|c|c|c|}
\hline
$P\left( x^2k \right)$ & 0.005 & 0.010 & 0.001 \\ \hline
$k$ & 3.841 & 6.635 & 10.828 \\ \hline
\end{tabular}
\end{center}
\end{question}
\begin{answer}{$p=0.9$；$\chi^2>10.828$，检测结果与患病有关。}
(1) $p=\dfrac{180}{200}=0.9$。(2) $\chi^2=\dfrac{1000(20\times20-180\times780)^2}{800\times200\times200\times800}>10.828$，故认为样本数据中超声波检测结果与患该疾病有关。
\end{answer}

\begin{question}[GPT生成]
从标有$1,2,3,4$的4张卡片中任取1张，取到偶数的概率为\choiceBox
\choices{$\dfrac14$}{$\dfrac13$}{$\dfrac12$}{$\dfrac34$}
\end{question}
\begin{answer}{C}
样本点共有 $4$ 个，取到偶数对应 $2,4$ 两个样本点，故概率为 $\dfrac24=\dfrac12$。选 C。
\end{answer}

\questionGroup{二、填空题}

% ----------------------------------------
\begin{question}[GPT生成]
抛掷一枚均匀硬币两次，恰好出现一次正面的概率为\blank{3cm}．
\end{question}
\begin{answer}{$\dfrac12$}
两次抛掷的样本点为 正正、正反、反正、反反，共 $4$ 个，其中恰好一次正面有 $2$ 个，故概率为 $\dfrac24=\dfrac12$。
\end{answer}

% ===== 9.2 统计 =====
\QuestionSetSection{统计}

\questionGroup{一、选择题}

% ----------------------------------------
\begin{question}[GPT生成]
一组数据$2,4,4,6,9$的中位数为\choiceBox
\choices{4}{5}{6}{9}
\end{question}
\begin{answer}{A}
数据已按从小到大排列，共 $5$ 个数，中位数是第 $3$ 个数，即 $4$。选 A。
\end{answer}

\questionGroup{三、解答题}

% ----------------------------------------
\begin{question}[GPT生成]
某班5名学生一次测试成绩分别为$70,80,80,90,100$，求这组数据的平均数和方差。
\end{question}
\begin{answer}{平均数 $84$，方差 $104$}
平均数为 $\bar x=\dfrac{70+80+80+90+100}{5}=84$。方差为 $\dfrac{(70-84)^2+(80-84)^2+(80-84)^2+(90-84)^2+(100-84)^2}{5}=\dfrac{196+16+16+36+256}{5}=104$。
\end{answer}

\QuestionSetSection*{本章综合训练}

\questionGroup{综合解答题}

% ----------------------------------------
\begin{question}[GPT生成]
袋中有大小相同的$3$个红球和$2$个蓝球，不放回任取$2$个。设$X$为取到红球的个数。（1）求两个都是红球的概率；（2）写出$X$的分布列；（3）求$E(X)$。
\end{question}
\solutionSpace{5cm}
\begin{answer}{(1) $\dfrac3{10}$；(2) $P(X=0)=\dfrac1{10},P(X=1)=\dfrac35,P(X=2)=\dfrac3{10}$；(3) $\dfrac65$}
总取法数为 $\mathrm{C}_5^2=10$，两个都是红球的取法为 $\mathrm{C}_3^2=3$，故概率为 $\dfrac3{10}$。$X=0$ 时取两蓝，概率 $\dfrac{\mathrm{C}_2^2}{10}=\dfrac1{10}$；$X=1$ 时取一红一蓝，概率 $\dfrac{\mathrm{C}_3^1\mathrm{C}_2^1}{10}=\dfrac35$；$X=2$ 的概率为 $\dfrac3{10}$。所以 $E(X)=0\cdot\dfrac1{10}+1\cdot\dfrac35+2\cdot\dfrac3{10}=\dfrac65$。
\end{answer}

% ----------------------------------------
\begin{question}[GPT生成]
一组数据$6,8,10,12,x$的平均数为$10$。（1）求$x$；（2）求这组数据的中位数；（3）求这组数据的方差，并说明方差反映了数据的什么特征。
\end{question}
\solutionSpace{5cm}
\begin{answer}{(1) $x=14$；(2) $10$；(3) 方差 $8$，反映数据离散程度}
由平均数为 $10$，得 $6+8+10+12+x=50$，所以 $x=14$。数据从小到大为 $6,8,10,12,14$，中位数为 $10$。方差为 $\dfrac{(6-10)^2+(8-10)^2+(10-10)^2+(12-10)^2+(14-10)^2}{5}=\dfrac{40}{5}=8$。方差越大，数据围绕平均数波动越明显。
\end{answer}

\printChapterAnswers
